\section{Conclusion}
\label{sec:conclusion}

We presented the first empirical mechanistic characterization of tool selection in \LLMs,
combining behavioral analysis, semantic probing, causal perturbation, and layer-wise
representational probing.  Our four experiments converge on a coherent picture.

At the behavioral level, tool selection is driven primarily by semantic description
content: cosine similarity predicts 40\% of selections (vs.\ 25\% chance), selection
stability across orderings is 90\%, and adversaially swapping descriptions changes
selection in 67\% of cases.  Positional bias exists in uncontrolled settings but is not
significant in controlled rotation experiments, indicating that high position-0 rates
observed in prior work partly reflect correlated description quality.  Tool names carry
an independent prior that can override description content for highly distinctive names.

At the mechanistic level, we uncover a clear dissociation within the \gpttwo transformer:
tool category is fully encoded from layer~1 (100\% probe accuracy), consistent with
lexical encoding in token embeddings.  In contrast, the identification of which specific
tool position is semantically correct emerges progressively, rising from 35\% at layer~1
to 68\% at layer~9, before declining.  This gradient establishes that semantic matching
is a non-trivial multi-layer computation concentrated in the middle-to-late transformer
layers.

\para{Future work.}  The most important next step is applying ACDC circuit discovery
\citep{conmy2023acdc} to identify the specific attention heads responsible for cross-tool
comparison in layers 7--9 of GPT-2-small.  Extending layer-wise probing to instruction-tuned
models (LLaMA-3.1-8B-Instruct, GPT-4.1-mini) would test whether tool-use fine-tuning
creates specialized selection circuits or reuses the same mid-layer comparison mechanism.
Sparse autoencoder (SAE) analysis \citep{sharkey2025open} could identify monosemantic
features that activate for tool-relevant concepts, enabling targeted activation steering
to correct selection errors without weight updates.  We release all datasets, code, and
experimental results to support replication and extension of this work.
